% Ejemplo ilustrativo (cifras de resultados hipotéticas). Sustitúyelo por el texto de tu investigación.

\guia{Un solo párrafo de 150 a 250 palabras, sin sangría, en este orden: objetivo, metodología, hallazgos, limitación y aportación. Debajo van las palabras clave (4 a 6, en minúsculas, separadas por punto y coma).}

El presente estudio tiene como objetivo analizar la relación entre la volatilidad del tipo de cambio peso-dólar y el volumen de negociación de contratos de futuros del dólar en el Mercado Mexicano de Derivados (MexDer) durante el período 2009–2022. Para ello se empleó una metodología cuantitativa de alcance correlacional-explicativo, con series mensuales del Banco de México, el MexDer y el CME Group; la volatilidad cambiaria se estimó con un modelo GARCH(1,1) y su efecto sobre el volumen negociado se evaluó mediante regresiones con errores estándar robustos (Newey-West) para dos subperíodos. Los resultados indican que, entre 2009 y 2015, un aumento de un punto porcentual en la volatilidad condicional se asoció con un incremento de 4.2\% en el volumen de futuros ($\beta = 0.042$, $p < 0.01$); sin embargo, entre 2016 y 2022 el efecto dejó de ser significativo ($\beta = 0.011$, $p = 0.21$), mientras que el volumen de futuros del peso en el CME mostró una relación negativa con la actividad del MexDer ($\beta = -0.351$, $p < 0.01$). El volumen anual de futuros en el MexDer cayó alrededor de 86\% en el período. Como principal limitación, no se dispone de información mensual sobre la cobertura cambiaria pactada en el mercado extrabursátil (OTC) local, que podría explicar parte de la sustitución observada. Esta investigación contribuye a la literatura sobre mercados de derivados en economías emergentes al mostrar que la contracción del MexDer no obedece a una menor demanda de cobertura, sino a la migración de la liquidez hacia mercados externos y extrabursátiles.

\hspace{1.27cm}\textit{Palabras clave:} mercado de derivados; futuros del dólar; volatilidad cambiaria; cobertura de riesgo; México\\
\hspace*{1.27cm}\textit{Clasificación JEL:} G13, G15, F31
